\hypertarget{supplementary-material}{%
\section{Supplementary material}\label{supplementary-material}}

\begin{itemize}
\tightlist
\item
  \textbf{S1. Four cohort definitions.} Four readings of ``first ICU
  admission at first hospitalization'', with counts and one-year
  mortality.
\item
  \textbf{S2. Reference cohort: complete results.} Nested
  cross-validation of all nine configurations, LASSO selection
  frequencies, the model-selection steps, the secondary held-out
  estimate, thresholds, decision curves, coefficients and SHAP.
\item
  \textbf{S3. Full baseline tables} of the reference and landmark
  cohorts.
\item
  \textbf{S4. Comparison with the benchmark's Table 1} and the
  mixed-unit diagnostic for the white-cell differential.
\item
  \textbf{S5. Hypertension definitions} and their associations with
  one-year death.
\item
  \textbf{S6. LASSO cross-validation curve and coefficient path.}
\item
  \textbf{S7. Exploratory single-split results} (all models, paired
  DeLong comparisons, thresholds, calibration, LASSO-only nesting).
\item
  \textbf{S8. Comorbidity-timing recode}: acute codes removed and
  changes in each comorbidity.
\item
  \textbf{S9. Implementation details.} Global random seed 42 and repeat
  seeds; LASSO implementation (λ grid of 100 values from
  \(\lambda_{\max} = \max_j \lvert\mathbf{x}_j^\top(\mathbf{y} - \bar y)\rvert / n\)
  down to \(10^{-4}\lambda_{\max}\); liblinear intercept feature set to
  1,000 so that the intercept penalty is negligible; binomial deviance
  as the criterion, whereas the benchmark's text names mean squared
  error but its figure shows binomial deviance); the benchmark's
  hyperparameters and their mapping to scikit-learn and xgboost; all
  tuning grids.
\item
  \textbf{S10. Benchmark-style versus within-fold preprocessing}: LASSO
  selection, validation AUROC and SHAP.
\item
  \textbf{S11. Code availability and reproducibility}: pipeline scripts,
  seeds and the frozen-model specification. All code, aggregate results
  and this manuscript's sources are available at
  \url{https://github.com/hongfangyu123/BMI5101-DN-Mortality-P16};
  patient-level MIMIC-IV data are not included and require credentialed
  PhysioNet access.
\item
  \textbf{S12. Exploratory extensions: complete results.} All 14 models
  (LR, RF, XGBoost and stacking with static, trajectory, burden and
  embedding inputs), all 45 paired comparisons, pooled out-of-fold
  calibration, decision-curve comparisons and trajectory-feature
  selection frequencies. The extension experiments were run on a
  separate machine; rebuilding their input features from our own data
  reproduced the same patients, outcomes and missingness pattern, with a
  maximum absolute difference of \(2.8 \times 10^{-14}\).
\item
  \textbf{S13. Post-hoc diagnostics of the held-out AUROC.} Exploratory
  analyses that did not establish the cause of the difference between
  nested-CV and held-out AUROC. The frozen LR's apparent training AUROC
  was 0.817. Resampling the OOF predictions 10,000 times to the held-out
  composition (136 deaths, 376 survivors) gave a mean AUROC of 0.800 (SD
  0.021; 2.5th percentile 0.755), with 1.4\% of resamples at or below
  0.750; this ignores refitting and selection. The largest standardised
  mean difference between training and held-out sets across the 22
  selected variables was 0.11; similar marginal distributions do not
  exclude differences in the joint distribution.
\end{itemize}
